\documentclass[11pt,a4paper]{article}
% Preâmbulo comum. Não duplicar pacotes nos módulos.
\usepackage[utf8]{inputenc}
\usepackage[T1]{fontenc}
\usepackage[brazil]{babel}
\usepackage{amsmath,amssymb,amsthm,mathtools}
\usepackage{geometry}
\geometry{a4paper,margin=2.3cm}
\usepackage{microtype}
\usepackage{enumitem}
\usepackage{graphicx}
\usepackage{xcolor}
\usepackage{booktabs}
\usepackage{hyperref}
\hypersetup{colorlinks=true,linkcolor=blue!40!black,urlcolor=blue!50!black}

\setcounter{tocdepth}{2}
\numberwithin{equation}{section}

\theoremstyle{plain}
\newtheorem{teorema}{Teorema}[section]
\newtheorem{proposicao}[teorema]{Proposição}
\newtheorem{lema}[teorema]{Lema}
\newtheorem{corolario}[teorema]{Corolário}

\theoremstyle{definition}
\newtheorem{definicao}[teorema]{Definição}
\newtheorem{exemplo}[teorema]{Exemplo}
\newtheorem{exercicio}{Exercício}[section]

\theoremstyle{remark}
\newtheorem{observacao}[teorema]{Observação}
\newtheorem{contraexemplo}[teorema]{Contraexemplo}
\newtheorem{intuicao}[teorema]{Intuição}

\newcommand{\R}{\mathbb{R}}
\newcommand{\N}{\mathbb{N}}
\newcommand{\Z}{\mathbb{Z}}
\newcommand{\Q}{\mathbb{Q}}
\DeclareRobustCommand{\eps}{\varepsilon}
\DeclareRobustCommand{\vet}[1]{\vec{#1}}
\DeclareRobustCommand{\dif}{\mathrm{d}}
\DeclareRobustCommand{\abs}[1]{\left\lvert #1 \right\rvert}
\DeclareRobustCommand{\norm}[1]{\left\lVert #1 \right\rVert}
\DeclareMathOperator{\grad}{grad}
\DeclareMathOperator{\Div}{div}
\DeclareMathOperator{\rot}{rot}

\newcommand{\base}[1]{\par\medskip\noindent\textbf{Base para relembrar.} #1\par\medskip}
\newcommand{\cuidado}[1]{\par\medskip\noindent\textbf{Quando a hipótese falha.} #1\par\medskip}

\newcommand{\figura}[2]{%
  \begin{figure}[htbp]
  \centering
  \IfFileExists{#1}{%
    \includegraphics[width=0.84\textwidth]{#1}%
  }{%
    \fbox{\parbox{0.8\textwidth}{\centering Figura ainda não gerada:\\ \texttt{#1}}}%
  }
  \caption{#2}
  \end{figure}%
}

\newenvironment{esqueleto}{%
  \par\medskip\noindent\textbf{Esqueleto da demonstração.}\par\nopagebreak
  \begin{enumerate}[leftmargin=1.4em]
}{\end{enumerate}\medskip}

\title{Gabarito --- Trigonometria}
\author{}
\date{}
\begin{document}
\maketitle
\tableofcontents

A numeração segue a ordem dos exercícios do texto.

\section*{Exercício 1. [Básico --- valores em outro quadrante]}

\textbf{Enunciado.} Calcule $\sin(5\pi/6)$, $\cos(5\pi/6)$ e $\tan(5\pi/6)$.

\textbf{Solução.}
O ângulo $5\pi/6$ está em $(\pi/2,\pi)$, e $5\pi/6=\pi-\pi/6$. As reduções dão
\[
\sin\frac{5\pi}{6}=\sin\Bigl(\pi-\frac{\pi}{6}\Bigr)=\sin\frac{\pi}{6}=\frac12,
\]
\[
\cos\frac{5\pi}{6}=-\cos\frac{\pi}{6}=-\frac{\sqrt{3}}{2}.
\]
O cosseno não é zero, então a tangente existe:
\[
\tan\frac{5\pi}{6}
=\frac{1/2}{-\sqrt{3}/2}
=-\frac{1}{\sqrt{3}}
=-\frac{\sqrt{3}}{3}.
\]
A última igualdade vale porque $\sqrt{3}/3$ e $1/\sqrt{3}$ são positivos e têm o mesmo quadrado $1/3$.

\section*{Exercício 2. [Básico --- equação do cosseno]}

\textbf{Enunciado.} Resolva $\cos x=\sqrt{2}/2$, com a solução geral.

\textbf{Solução.}
O valor notável $\cos(\pi/4)=\sqrt{2}/2$ fornece uma solução particular. A solução geral do cosseno é
\[
x=\pm\frac{\pi}{4}+2k\pi,\qquad k\in\Z.
\]
A ida: $\cos$ é par e tem período $2\pi$, logo $\cos(\pi/4+2k\pi)=\cos(-\pi/4+2k\pi)=\cos(\pi/4)$. A volta é o teorema geral: se $\cos x=\cos(\pi/4)$, então $x=\pi/4+2k\pi$ ou $x=-\pi/4+2k\pi$.

\section*{Exercício 3. [Básico --- arco metade]}

\textbf{Enunciado.} Calcule $\sin(\pi/12)$ e $\cos(\pi/12)$ com sinal pelo quadrante.

\textbf{Solução.}
Tem-se $\pi/12=(1/2)\cdot(\pi/6)$ e $\pi/12\in(0,\pi/2)$, logo seno e cosseno são positivos. As fórmulas de metade dão
\begin{align*}
\sin\frac{\pi}{12}
&=\sqrt{\frac{1-\cos(\pi/6)}{2}}
=\sqrt{\frac{1-\sqrt{3}/2}{2}}
=\sqrt{\frac{2-\sqrt{3}}{4}}
=\frac{\sqrt{2-\sqrt{3}}}{2},\\[0.4em]
\cos\frac{\pi}{12}
&=\sqrt{\frac{1+\cos(\pi/6)}{2}}
=\sqrt{\frac{2+\sqrt{3}}{4}}
=\frac{\sqrt{2+\sqrt{3}}}{2}.
\end{align*}
Esses valores coincidem com $(\sqrt{6}-\sqrt{2})/4$ e $(\sqrt{6}+\sqrt{2})/4$. De fato, ambos os candidatos a seno são positivos, e
\begin{align*}
\Bigl(\frac{\sqrt{6}-\sqrt{2}}{4}\Bigr)^2
&=\frac{6+2-2\sqrt{12}}{16}
=\frac{8-4\sqrt{3}}{16}
=\frac{2-\sqrt{3}}{4}
=\Bigl(\frac{\sqrt{2-\sqrt{3}}}{2}\Bigr)^2.
\end{align*}
A unicidade da raiz não negativa identifica os dois números. O mesmo cálculo com o sinal interno trocado identifica os cossenos:
\[
\Bigl(\frac{\sqrt{6}+\sqrt{2}}{4}\Bigr)^2=\frac{8+4\sqrt{3}}{16}=\frac{2+\sqrt{3}}{4}.
\]

\section*{Exercício 4. [Básico --- domínio trigonométrico]}

\textbf{Enunciado.} Determine o domínio de $\tan x+\sec x$ e o de $\csc(2x)$.

\textbf{Solução.}
A soma $\tan x+\sec x$ está definida exatamente quando as duas parcelas estão definidas. As duas exigem $\cos x\neq 0$, isto é $x\neq \pi/2+k\pi$ para todo $k\in\Z$. Esse é o domínio.

A função $\csc(2x)$ está definida exatamente quando $\sin(2x)\neq 0$. Isso ocorre se e somente se $2x\neq n\pi$ para todo $n\in\Z$, isto é $x\neq n\pi/2$. O domínio é $\R\setminus\{n\pi/2:n\in\Z\}$.

\section*{Exercício 5. [Intermediário --- soma de cossenos]}

\textbf{Enunciado.} Prove $\cos a+\cos b=2\cos\frac{a+b}{2}\cos\frac{a-b}{2}$.

\textbf{Solução.}
As fórmulas de adição e de subtração dão
\begin{align*}
\cos(x+y)&=\cos x\cos y-\sin x\sin y,\\
\cos(x-y)&=\cos x\cos y+\sin x\sin y.
\end{align*}
Somando, os termos com seno cancelam-se e
\[
\cos(x+y)+\cos(x-y)=2\cos x\cos y.
\]
Ponhamos $a=x+y$ e $b=x-y$. Então $2x=a+b$ e $2y=a-b$, logo $x=(a+b)/2$ e $y=(a-b)/2$. Substituindo,
\[
\cos a+\cos b=2\cos\frac{a+b}{2}\cos\frac{a-b}{2}.
\]

\section*{Exercício 6. [Intermediário --- lei dos cossenos numérica]}

\textbf{Enunciado.} Dois lados medem $2$ e $3$ e o ângulo entre eles mede $\pi/3$. Calcule o terceiro lado e o cosseno do ângulo oposto ao lado $2$.

\textbf{Solução.}
Sejam $a=2$, $b=3$ e $\gamma=\pi/3$ o ângulo entre eles. O lado $c$ oposto a $\gamma$ satisfaz
\[
c^2=2^2+3^2-2\cdot 2\cdot 3\cdot\cos\frac{\pi}{3}
=4+9-12\cdot\frac12
=13-6=7.
\]
Como $c>0$, $c=\sqrt{7}$.

Seja $\alpha$ o ângulo oposto ao lado $2$. A lei dos cossenos, agora com o lado $2$ isolado, usa os outros dois lados $3$ e $\sqrt{7}$:
\[
2^2=3^2+(\sqrt{7})^2-2\cdot 3\cdot\sqrt{7}\cdot\cos\alpha,
\]
isto é $4=9+7-6\sqrt{7}\cos\alpha=16-6\sqrt{7}\cos\alpha$. Então
\[
6\sqrt{7}\cos\alpha=12,\qquad \cos\alpha=\frac{2}{\sqrt{7}}=\frac{2\sqrt{7}}{7}.
\]
O número é positivo e menor do que $1$, porque $2\sqrt{7}<2\cdot 3=6<7$. O ângulo $\alpha$, estando em $(0,\pi)$, fica em $(0,\pi/2)$.

\section*{Exercício 7. [Intermediário --- tangente da soma]}

\textbf{Enunciado.} Prove a fórmula de $\tan(a+b)$ sob as hipóteses $\cos a\cos b\neq 0$ e $\cos(a+b)\neq 0$.

\textbf{Solução.}
A hipótese $\cos(a+b)\neq 0$ é exatamente a condição para $\tan(a+b)$ estar definida, e permite escrever
\[
\tan(a+b)=\frac{\sin(a+b)}{\cos(a+b)}
=\frac{\sin a\cos b+\cos a\sin b}{\cos a\cos b-\sin a\sin b}.
\]
A hipótese $\cos a\cos b\neq 0$ diz que $\cos a\neq 0$ e $\cos b\neq 0$, logo $\tan a$ e $\tan b$ existem, e permite dividir o numerador e o denominador por $\cos a\cos b$:
\begin{align*}
\frac{\sin a\cos b}{\cos a\cos b}+\frac{\cos a\sin b}{\cos a\cos b}
&=\tan a+\tan b,\\[0.3em]
\frac{\cos a\cos b}{\cos a\cos b}-\frac{\sin a\sin b}{\cos a\cos b}
&=1-\tan a\tan b.
\end{align*}
Portanto
\[
\tan(a+b)=\frac{\tan a+\tan b}{1-\tan a\tan b}.
\]
Se $\cos a\cos b=0$, a divisão não está definida e pelo menos uma das tangentes do segundo membro não existe. Se $\cos(a+b)=0$, o primeiro membro não existe; é o que ocorre para $a=b=\pi/4$, em que $\tan a$ e $\tan b$ existem, o denominador $1-1\cdot 1$ é zero e $a+b=\pi/2$.

\section*{Exercício 8. [Avançado --- soma de arctangentes]}

\textbf{Enunciado.} Prove que $\arctan 1+\arctan 2+\arctan 3=\pi$, separando $\arctan 2+\arctan 3$.

\textbf{Solução.}
Sejam $\alpha=\arctan 2$ e $\beta=\arctan 3$. Ambos estão em $(0,\pi/2)$. Como $\tan$ é estritamente crescente nesse intervalo e $\tan(\pi/4)=1<2=\tan\alpha$, segue $\alpha>\pi/4$. Do mesmo modo $\beta>\pi/4$. Logo $\alpha+\beta>\pi/2$. Cada um é menor do que $\pi/2$, então $\alpha+\beta<\pi$ e $\alpha+\beta\in(\pi/2,\pi)$. Nesse intervalo o cosseno não se anula: $\cos\theta=0$ só em $\pi/2+k\pi$, e nenhum desses pontos está em $(\pi/2,\pi)$.

Como $\alpha\in(0,\pi/2)$ e $\tan\alpha=2$, tem-se $\cos\alpha>0$ e $\sin\alpha=2\cos\alpha$. A equação fundamental dá $4\cos^2\alpha+\cos^2\alpha=1$, logo $\cos\alpha=1/\sqrt{5}$ e $\sin\alpha=2/\sqrt{5}$. Do mesmo modo $\cos\beta=1/\sqrt{10}$ e $\sin\beta=3/\sqrt{10}$. As fórmulas de adição produzem
\begin{align*}
\sin(\alpha+\beta)
&=\frac{2}{\sqrt{5}}\cdot\frac{1}{\sqrt{10}}+\frac{1}{\sqrt{5}}\cdot\frac{3}{\sqrt{10}}
=\frac{5}{\sqrt{50}}
=\frac{1}{\sqrt{2}},\\
\cos(\alpha+\beta)
&=\frac{1}{\sqrt{5}}\cdot\frac{1}{\sqrt{10}}-\frac{2}{\sqrt{5}}\cdot\frac{3}{\sqrt{10}}
=\frac{-5}{\sqrt{50}}
=-\frac{1}{\sqrt{2}},
\end{align*}
porque $\sqrt{50}=5\sqrt{2}$ e $5/(5\sqrt{2})=1/\sqrt{2}$. Portanto $\tan(\alpha+\beta)=-1$. O ângulo $3\pi/4$ está em $(\pi/2,\pi)$ e satisfaz $\sin(3\pi/4)=\sqrt{2}/2=1/\sqrt{2}$ e $\cos(3\pi/4)=-\sqrt{2}/2$. A parametrização do círculo é injetora em intervalos de comprimento menor do que $2\pi$, e $(\pi/2,\pi)$ tem comprimento menor do que $2\pi$, então há um único ângulo nesse intervalo com esse seno e esse cosseno. Logo $\alpha+\beta=3\pi/4$.

O valor $-\pi/4$ também tem tangente $-1$, mas tem seno negativo e não está em $(\pi/2,\pi)$. A soma de dois ângulos maiores do que $\pi/4$ não pode ser $-\pi/4$.

Por fim $\arctan 1=\pi/4$, porque $\tan(\pi/4)=1$ e $\pi/4\in(-\pi/2,\pi/2)$. Portanto
\[
\arctan 1+\arctan 2+\arctan 3=\frac{\pi}{4}+\frac{3\pi}{4}=\pi.
\]

\section*{Exercício 9. [Avançado --- equação com soma para produto]}

\textbf{Enunciado.} Resolva $\sin(5x)-\sin x=0$ como união de duas famílias, sem raízes extras.

\textbf{Solução.}
A identidade da diferença de senos, com $a=5x$ e $b=x$, dá
\[
\sin(5x)-\sin x
=2\cos\frac{5x+x}{2}\sin\frac{5x-x}{2}
=2\cos(3x)\sin(2x).
\]
A equação $2\cos(3x)\sin(2x)=0$ equivale a $\cos(3x)=0$ ou $\sin(2x)=0$, porque um produto de reais é zero se e somente se um dos fatores é zero. Não houve divisão, então nenhuma raiz foi introduzida nem perdida: os dois lados são idênticos como funções.

A equação $\cos(3x)=0$ equivale a $3x=\pi/2+k\pi$ para algum $k\in\Z$, isto é
\[
x=\frac{\pi}{6}+\frac{k\pi}{3},\qquad k\in\Z.
\]
A equação $\sin(2x)=0$ equivale a $2x=n\pi$ para algum $n\in\Z$, isto é
\[
x=\frac{n\pi}{2},\qquad n\in\Z.
\]
A solução geral é a união dessas duas famílias. Há sobreposição: por exemplo $x=\pi/2$ pertence à segunda com $n=1$ e à primeira com $k=1$, pois $\pi/6+\pi/3=\pi/2$. A união não exige que as famílias sejam disjuntas. Um ângulo que anula os dois fatores continua sendo raiz da equação original, porque o produto é zero.

\end{document}
